\subsection{Widerstandsberechnung für nicht-invertierenden Verstärker}
% Alt: Aufgabe 2
\Aufgabe{
	\label{Aufg2}
	Gegeben ist die abgebildete Schaltung mit einem nichtinvertierenden
	Verstärker. Dabei ist der Widerstand $R_\mathrm{2} = 1\,\mathrm{k \Omega}$. Was für Widerstandswerte
	muss $R_\mathrm{1}$ aufweisen, damit die Verstärkung der gegebenen Schaltung 10, 50 und 100 beträgt?

	\begin{figure}[H]
		\centering
		\input{Tikz/src/FigWiderstandNIV.tex}\alttext{Das Bild zeigt ein elektrisches Schaltbild mit einem Operationsverstärker, der als Spannungsfolger geschaltet ist. Links ist eine Eingangsspannung \(U_\mathrm{E}\) mit einem Pfeil nach unten eingezeichnet. Der Plus-Eingang des Operationsverstärkers ist mit der Eingangsspannung verbunden. Der Minus-Eingang ist über einen Rückkopplungszweig direkt mit dem Ausgang des Operationsverstärkers verbunden. Am Ausgang des Operationsverstärkers befinden sich zwei in Serie geschaltete Widerstände \(R_2\) und \(R_1\), die zum Massepunkt führen, welcher unten in der Mitte geerdet ist. Rechts befindet sich die Ausgangsspannung \(U_\mathrm{A}\) mit einem Pfeil nach unten.}
		\label{fig:FigWiderstandNIV}
	\end{figure}

}

\Loesung{
	Gemäß Formel nichtinvertierender Verstärker:
	\begin{align*}
		V = 1 + \frac{R_\mathrm{1}}{R_\mathrm{2}}
	\end{align*}
	Formel nach $R_1$ umstellen:
	\begin{align*}
		R_1 = R_2 \cdot (V - 1)
	\end{align*}
	\begin{itemize}
		\item $V = 10$
		      \begin{align*}
			      R_1 = \mathrm{1\,k\Omega \cdot (10 - 1) = 9\,k\Omega}
		      \end{align*}
		\item $V = 50$
		      \begin{align*}
			      R_1 = \mathrm{1\,k\Omega \cdot (50 - 1) = 49\,k\Omega}
		      \end{align*}
		\item $V = 100$
		      \begin{align*}
			      R_1 = \mathrm{1\,k\Omega \cdot (100 - 1) = 99\,k\Omega}
		      \end{align*}
	\end{itemize}

	Siehe: Abschnitt \ref{fig: Verschaltung Verstaerker} und Tabelle \ref{tab:Grundschaltungen}
}